% ! TEX root = ../bcd-19-20.pdf

\chapter{Parțial 2018-2019}

\section{Numărul 1}

\indent\indent 1. Studiați natura seriei:
$ \ds\sum_{n \geq 1} \dfrac{\alpha^n}{2n^2 + 1} $.

\vspace{1cm}

2. Studiați convergența absolută și convergența seriei
$ \ds\sum_{n \geq 1} (-1)^n \dfrac{(2n)!}{n^n} $.

\vspace{1cm}

3. Studiați convergența simplă și uniformă a șirului de funcții:
\[
  f_n : [0, 2] \to \RR, \quad f_n(x) = \dfrac{nx^3}{x^2 + n}.
\]

\vspace{1cm}

4. Studiați convergența uniformă pentru seria de funcții
$ \ds\sum_{n \geq 1} \dfrac{xe^{-nx}}{\sqrt{n}}, x \geq 0 $.

\vspace{1cm}

5. Aproximați cu o eroare de maximum $ \varepsilon = 10^{-3} $ soluția
reală a ecuației $ x^3 + 5x - 1 = 0 $.

%%%%%%%%%%%%%%%%%%%%%%%%%%%%%%%%%%%%%%%%%%%%%%%%%%%%%%%%%%%%%%%%%%%%%%

\section{Numărul 2}

\indent\indent 1. Studiați natura seriei:
$ \ds\sum_{n \geq 1} \left( \dfrac{n}{n+a} \right)^{n^2}, a > 0 $.

\vspace{1cm}

2. Studiați convergența absolută și convergența seriei
$ \ds\sum_{n \geq 1} \dfrac{(-1)^n}{2^{2n} \cdot n!} $.

\vspace{1cm}

3. Studiați convergența simplă și convergența uniformă pentru șirul de funcții:
\[
  f_n : [0, 1] \to \RR, \quad f_n(x) = \dfrac{nx^n + 1}{nx^n + 2}.
\]

\vspace{1cm}

4. Studiați convergența uniformă a seriei de funcții:
$ \ds\sum_{n \geq 1} \dfrac{x}{1 + n^4 x^2}, x \geq 0 $.

\vspace{1cm}

5. Aproximați cu o eroare de maximum $ \varepsilon = 10^{-3} $ soluția reală
a ecuației: $ x^3 + 3x - 1 = 0 $.

%%%%%%%%%%%%%%%%%%%%%%%%%%%%%%%%%%%%%%%%%%%%%%%%%%%%%%%%%%%%%%%%%%%%%%

\section{Numărul 3}

\indent\indent 1. Studiați natura seriei:
$ \ds\sum_{n \geq 1} \dfrac{n\sqrt{n}}{\left(a + \dfrac{1}{n} \right)^n}, a > 0 $.

\vspace{1cm}

2. Studiați convergența absolută a seriei: $ \ds\sum_{n \geq 1} (-1)^n \dfrac{(n!)^2}{(2n)!} $.

\vspace{1cm}

3. Studiați convergența simplă și uniformă pentru șirul de funcții:
\[
  f_n : [1, \infty) \to \RR, \quad f_n(x) = \dfrac{n}{x^3 + nx}.
\]

\vspace{1cm}

4. Studiați convergența uniformă a seriei de funcții:
$ \ds\sum_{n \geq 1} \dfrac{(-1)^n x^2}{1 + n^3 x^4}, x \in \RR $.

\vspace{1cm}

5. Să se aproximeze cu o eroare de maximum $ \varepsilon = 10^{-3} $ soluția
reală a ecuației $ x^3 + 8x - 1 = 0 $.



%%% Local Variables:
%%% mode: latex
%%% TeX-master: "../bcd-19-20"
%%% End:
